% ---- pg-035 （书页 23） ----
见文献[30]中的讨论.

有许多灵巧的方法可以在二维空间中为多元数据作图. 我们已介绍其中几种, 利用先进的计算机作图是可能的而且几乎一定会取得更大的进展.

\section{距离}  % 原书 1.5
\phantomsection\label{sec:1.5}

尽管距离最初显得很难, 但大多数多元方法都基于简单的距离概念. 直线或欧几里德距离是我们所熟悉的. 如果我们考虑平面上的点 $P(x_1, x_2)$, 从 $P$ 点到原点 $O(0,0)$ 的直线距离 $d(O, P)$, 根据勾股定理, 是
\phantomsection\label{eq:1-9}
\begin{equation}
d(O, P) = \sqrt{x_1^2 + x_2^2} \tag{1-9}
\end{equation}
这种情况用\hyperref[fig:1.19]{图 1.19} 说明. 一般的, 若点 $P$ 有 $p$ 个坐标, 则 $P = (x_1, x_2, \cdots, x_p)$, $P$ 到原点 $O = (0, 0, \cdots, 0)$ 的直线距离为
\phantomsection\label{eq:1-10}
\begin{equation}
d(O, P) = \sqrt{x_1^2 + x_2^2 + \cdots + x_p^2} \tag{1-10}
\end{equation}
(参见第 2 章.)离原点有常数距离的平方(例如 $c^2$)的所有点 $(x_1, x_2, \cdots, x_p)$ 满足方程
\phantomsection\label{eq:1-11}
\begin{equation}
d^2(O, P) = x_1^2 + x_2^2 + \cdots + x_p^2 = c^2 \tag{1-11}
\end{equation}
由于这是一个超球面的方程(当 $p=2$ 时是个圆), 与原点等距的点都在超球面上.

坐标为 $P = (x_1, x_2, \cdots, x_p)$ 和 $Q = (y_1, y_2, \cdots, y_p)$ 的任意两个点 $P$ 和 $Q$ 的直线距离为
\phantomsection\label{eq:1-12}
\begin{equation}
d(P, Q) = \sqrt{(x_1 - y_1)^2 + (x_2 - y_2)^2 + \cdots + (x_p - y_p)^2} \tag{1-12}
\end{equation}

\begin{mvfigure}
\mvinc{pg-035}{0.604,0.301,0.881,0.3875}
\phantomsection\label{fig:1.19}\mvcap{图 1.19\quad 由勾股定理给出的距离}
\end{mvfigure}

直线或欧几里德距离对大多数的统计目的是不合用的. 这是因为每个坐标对于计算欧几里德距离起着相等的作用. 当坐标被用来表示不同随机波动的测量值时, 经常需要对那些可变性大的坐标加权大大少于对没有很大变化性的坐标的加权, 这就需要一种距离的不同的度量.

我们现在的目的是导出一个“统计的”距离. 这是由于考虑到变差的不同以及相关性的存在. 由于我们的选择取决于样本方差及协方差, 在这一点上我们用术语统计距离来将它和通常的欧几里德里距离区分开. 统计距离正是多元分析的基础.

开始, 我们把作 $p$ 维散布图的观测值集合看做是固定的. 根据这些我们将构造从原点到一点 $P = (x_1, x_2, \cdots, x_p)$ 的距离的度量. 在我们的论证中, $P$ 的坐标 $(x_1, x_2, \cdots, x_p)$ 可以变化以产生该点的不同位置. 然而, 确定距离的数据仍然保持固定.

\begin{mvfigure}
\mvinc{pg-035}{0.142,0.698,0.410,0.8195}
\phantomsection\label{fig:1.20}\mvcap{图 1.20\quad 在 $x_1$ 方向上比在 $x_2$ 方向上变化性更大的散布图}
\end{mvfigure}

为举例说明, 假设我们有两个变量的 $n$ 对测量值, 称两个变量为 $x_1$ 和 $x_2$, 并假定变量 $x_1$ 的测量值关于 $x_2$ 的测量值独立地变化.\footnote{在这方面,“独立地”意指 $x_2$ 测量结果不能以任何的精确度预言来自 $x_1$ 的测量结果, 反之亦然.} 此外, 假设 $x_1$ 测量结果的可变性大于 $x_2$ 的测量结果的可变性. 数据的散布图看上去有点像\hyperref[fig:1.20]{图 1.20} 中的图片.

观察\hyperref[fig:1.20]{图 1.20}, 我们发现, 从原点出发向 $x_1$ 的方向给定的偏差的值不像从原点向 $x_2$ 方向等距离的值那样“惊人”或“不正常”. 这是因为 $x_1$ 方向内在的可变性大
